\chapter{State of the Art}

\section{Introduction and Scope of the Chapter}

This chapter maps the technical background for a reinforcement-learning-based cybersecurity defender that operates on flow-level network observations. The scope is deliberately broader than a comparison of binary classifiers. The thesis combines network intrusion detection, flow-feature representation, public benchmark datasets, supervised machine learning baselines, a Gymnasium-style dataset-as-environment formulation, QRDQN, leakage-aware preprocessing, stricter validation checks, data-efficiency experiments, and planned or preferred validation on independently captured laboratory traffic.

The chapter therefore follows a funnel. It first reviews network intrusion detection and traffic representation, then situates CICIDS2017 among public datasets, then discusses supervised and deep learning baselines, and finally places the RL formulation within both reinforcement-learning theory and prior RL/DRL work for intrusion detection and cyber defense. The final sections focus on methodological risks, evaluation protocols, training-scale evidence, and the research gap addressed by this thesis.

The central claim is cautious. Reinforcement learning for intrusion detection already exists, and public NIDS benchmarks are widely used \parencite{Yang2024DRLNIDSSurvey,Kheddar2025RLIDSReview,Gueriani2024DRLIoTIDSSurvey}. This thesis does not claim to introduce the first RL IDS, to prove QRDQN as the NIDS state of the art, or to demonstrate production-ready inline prevention. It studies a reproducible, offline, flow-level \texttt{PERMIT}/\texttt{BLOCK} decision pipeline and evaluates how that formulation behaves under controlled internal benchmarks, baseline comparison, and progressively stricter validation.

\section{Intrusion Detection and Network Security Monitoring}

Intrusion Detection Systems monitor host or network activity for evidence of attacks, misuse, or policy violations \parencite{ScarfoneMell2007}. A Network Intrusion Detection System focuses on traffic crossing network links or segments. It may inspect packets, flows, or derived events, and its output usually feeds a broader monitoring workflow such as alert triage, incident response, or later forensic analysis. This is different from an Intrusion Prevention System, which sits inline and can actively block, modify, or rate-limit traffic \parencite{ScarfoneMell2007}.

Classical IDS taxonomies distinguish signature-based, anomaly-based, specification-based, and hybrid approaches \parencite{ScarfoneMell2007}. Signature-based detection matches known malicious patterns and is useful for well-understood threats, but it depends on rule maintenance and is weak against novel or highly variable behavior. Anomaly-based detection models expected behavior and flags deviations; this makes it attractive for machine learning, but it is vulnerable to false positives, changing benign traffic, weak labels, and distribution shift \parencite{SommerPaxson2010ClosedWorld,Ring2019DatasetSurvey}. Specification-based methods encode allowed behavior explicitly, while hybrid systems combine rule-based and learned signals.

Modern NIDS work is therefore not only an algorithmic problem. It is also a measurement and deployment problem: what is observed, how it is represented, which errors matter, how alerts are consumed, and whether benchmark conditions resemble the intended environment. For this thesis, the implemented system should be read as an offline NIDS-like decision model over extracted flow rows. It produces \texttt{PERMIT}/\texttt{BLOCK} predictions in experiments, but it does not perform packet dropping, firewall updates, SDN rule installation, or production IPS enforcement.

\section{Network Traffic Representation}

Network traffic can be represented as packets, payload bytes, sessions, connections, or flows. Packet and payload representations preserve more detail, including header order, timing, and sometimes application-layer content, but they are costly to store and process at scale. Payload inspection also raises privacy concerns and becomes less available when traffic is encrypted. Flow representation aggregates packets into connection-like records and summarizes behavior through statistics such as duration, bytes, packets, inter-arrival times, packet-length distributions, TCP flags, and active or idle periods \parencite{Sperotto2010FlowIDS,Umer2017FlowIDS}.

Flow-based NIDS is attractive because it turns traffic into tabular observations that are compatible with classical machine learning, deep learning, and RL-style environments \parencite{Sperotto2010FlowIDS,LiuLang2019IDSSurvey}. CICIDS2017 follows this pattern: packet captures are processed into bidirectional flow CSVs using CICFlowMeter-style features \parencite{Sharafaldin2018CICIDS2017,Lashkari2017CICFlowMeter,CICIDS2017Web}. This is the representation used by the thesis. Each row is mapped into a fixed canonical vector and then presented to the learning system as one observation.

The convenience of flow representation comes with limits. A flow record cannot preserve all payload semantics or full packet-sequence behavior, so attacks whose evidence is content-level or highly sequential may be invisible or only indirectly represented \parencite{Sperotto2010FlowIDS,Umer2017FlowIDS}. Flow features can also encode dataset artifacts. IP addresses, absolute timestamps, Flow IDs, endpoint identity, capture-day structure, and ports used as scenario proxies can leak labels rather than represent attack behavior \parencite{Arp2020DosDontsMLSecurity,Kaufman2011LeakageDataMining}.

The current project uses a canonical flow schema with 76 feature values and a 76-value missingness mask, producing 152-dimensional observations. The mask records whether each canonical value was present and valid, rather than silently hiding missing or imputed fields. This design is not presented as a universal optimum for NIDS. Its purpose is reproducibility: the same feature contract is used for CICIDS2017 and for later flow inputs extracted from laboratory traffic.

\section{Public Datasets for Network Intrusion Detection}

Public NIDS datasets enable shared experimentation, reproducibility, and comparison across papers \parencite{Ring2019DatasetSurvey,Thakkar2020DatasetsReview}. They also constrain conclusions. A dataset reflects its traffic generator, topology, capture process, feature extractor, labels, attack scripts, and release format. Performance on a public benchmark is therefore evidence about a controlled task, not direct evidence of deployment readiness \parencite{Apruzzese2022CrossEvaluationNIDS,Layeghy2023CrossDomainNIDS}.

KDDCup99 and NSL-KDD are historically important because they shaped early IDS evaluation \parencite{KDDCup1999,Tavallaee2009NSLKDD}. NSL-KDD reduced some redundancy problems in KDDCup99, but both datasets are now limited as evidence for modern flow-based defense because their traffic, attacks, and feature schema are old \parencite{Tavallaee2009NSLKDD,Ring2019DatasetSurvey}. In this repository, NSL-KDD is therefore kept as historical benchmarking context, not as the final path for the current defender.

UNSW-NB15 is a more recent cyber-range dataset and remains useful for discussing modern benchmark design \parencite{MoustafaSlay2015UNSWNB15,Moustafa2016UNSWEval}. CICIDS2017 and CSE-CIC-IDS2018 are widely used CIC-family datasets for flow-based NIDS research \parencite{Sharafaldin2018CICIDS2017,CSECICIDS2018AWS}. Bot-IoT and ToN\_IoT broaden the landscape toward IoT and IIoT traffic \parencite{Koroniotis2019BotIoT,Alsaedi2020ToNIoT}. None of these datasets should be treated as a complete proxy for all real networks. Their value is highest when they are used with transparent preprocessing, strong baselines, and validation protocols that test distribution shift rather than only random matched-distribution performance.

\section{CICIDS2017 as the Main Internal Benchmark}

CICIDS2017 is the main internal benchmark in this thesis because it provides labeled benign and attack traffic, raw PCAPs, and flow CSVs generated with CICFlowMeter-style features \parencite{Sharafaldin2018CICIDS2017,Lashkari2017CICFlowMeter,CICIDS2017Web}. It was designed to improve on older IDS datasets by using contemporary traffic profiles and multiple attack scenarios \parencite{Sharafaldin2018CICIDS2017,Sharafaldin2019Detailed}. These properties make it suitable for controlled experimentation with flow-level NIDS models.

The thesis uses CICIDS2017 as an internal benchmark, not as proof that the defender is ready for deployment. This distinction is essential. CICIDS2017 traffic was generated in a controlled environment, and later studies have identified problems in its labels, flow extraction, feature calculation, duplicated or inconsistent flows, and common uses of the summarized CSV release \parencite{Engelen2021CICIDSIssues,Lanvin2023Errors,Liu2022ErrorPrevalenceNIDS,Rosay2022CICIDS2017Analysis,Dube2024FaultyCICIDS2017}. These critiques do not make the dataset useless, but they require careful documentation of the exact version, cleaning steps, feature exclusions, and split protocol.

The current repository responds to these concerns through a curated tracked dataset copy, a separate untracked raw reference copy, and loader-level cleaning in \texttt{src/load\_cicids2017.py}. The adapter performs numeric coercion, invalid-value handling, canonical mapping, and anti-leakage exclusion of fields that should not be model features. The evaluation code also supports random splits, CSV/day-style splits, and leave-one-exact-CSV-out validation. These steps reduce common risks, but they do not eliminate all benchmark artifacts.

For that reason, CICIDS2017 results should be presented as part of a validation ladder. Random split performance is useful for checking whether the model can learn the internal benchmark. Harder day/file splits, shuffled-label checks, leave-one-exact-CSV-out validation, and external lab-flow inference test different failure modes. A high score on the easiest setting is not sufficient for a deployment claim.

\section{Supervised Machine Learning and Deep Learning for NIDS}

Most flow-based NIDS work is naturally framed as supervised learning: each traffic record is mapped to a benign, malicious, or attack-family label \parencite{LiuLang2019IDSSurvey,Ahmad2021}. Classical supervised methods include decision trees, Random Forests, SVMs, k-NN, logistic regression, Naive Bayes, and gradient-boosted ensembles. On tabular flow features, tree-based models remain strong baselines because they handle non-linear interactions, heterogeneous feature scales, and irrelevant features well \parencite{LiuLang2019IDSSurvey,Apruzzese2022CrossEvaluationNIDS}.

This matters directly for an RL-based thesis. If each observation is a labeled flow row and the reward is derived from classification correctness, a supervised classifier is the natural baseline. A QRDQN agent cannot be evaluated only against its own reward curve or against weak baselines. It must be compared with strong supervised methods under the same feature schema, split protocol, preprocessing, and metrics. The current repository includes a Random Forest baseline protocol aligned with the QRDQN splits, while the maintained results snapshot still marks the latest full metrics as pending.

Deep learning expands the design space through MLPs, sequence models, autoencoders, attention mechanisms, transformer-style models, and graph-based approaches \parencite{Xiao2021DLIDSSurvey,Asadullah2023MLIDSSurvey,Sayeeduddin2022AINIDS,Nguyen2022BERTFlowNIDS}. These approaches can report very high benchmark scores, especially on public datasets, but surveys also identify recurring weaknesses: unclear preprocessing, favorable random splits, missing class-imbalance analysis, limited external validation, and insufficient reproducibility \parencite{LiuLang2019IDSSurvey,Ring2019DatasetSurvey,Arp2020DosDontsMLSecurity}.

Feature standardisation is another important supervised-learning lesson. Different datasets and tools expose overlapping but non-identical feature sets, which makes cross-dataset comparison difficult. Sarhan et al. proposed a NetFlow-aligned feature mapping to improve generalisability and explainability across NIDS datasets \parencite{Sarhan2021}. This thesis does not adopt that exact standard, but it adopts the same motivation: a fixed canonical schema makes the model interface auditable and allows CICIDS2017 and laboratory flow inputs to be compared through one preprocessing contract.

\section{Reinforcement Learning Foundations}

Reinforcement learning studies how an agent selects actions in an environment to maximize return \parencite{SuttonBarto2018RL}. Its core concepts are states or observations, actions, rewards, policies, value functions, and transitions. In a Markov decision process, the current state summarizes the information needed for future decision-making, and actions influence both immediate reward and later states \parencite{SuttonBarto2018RL}. This sequential property is what distinguishes RL from ordinary supervised classification.

Classical Q-learning is a model-free, off-policy temporal-difference method that learns action values without requiring a model of the environment dynamics \parencite{WatkinsDayan1992QLearning}. This lineage is relevant because DQN and QRDQN inherit the value-based, off-policy framing. Experience replay, introduced earlier as a mechanism for reusing past experience, also became central to later deep Q-learning systems \parencite{Lin1992ReactiveAgents}.

Cybersecurity can contain genuine sequential RL problems. A defender may isolate a host, patch a service, change routing, or deploy a countermeasure, and those actions can affect later observations and attacker opportunities. However, a static labeled-flow dataset is a weaker setting. If one flow row is sampled, classified, and followed by another row that is not causally affected by the previous action, the problem resembles cost-sensitive classification or a contextual decision process more than full autonomous cyber defense \parencite{Levine2020OfflineRL,Fujimoto2019BCQ}. The thesis must keep this distinction explicit.

\section{Deep Q-Learning and Distributional Reinforcement Learning}

Deep Q-Learning replaces a tabular action-value function with a neural network, making value-based RL applicable to high-dimensional observations \parencite{Mnih2015DQN}. DQN popularized two stabilizing mechanisms: replay memory and a separate target network \parencite{Mnih2015DQN}. Later variants addressed known weaknesses. Double DQN reduces overestimation bias by decoupling action selection from target evaluation \parencite{VanHasselt2016DoubleDQN}. Dueling DQN separates state value and action advantage estimation \parencite{Wang2016DuelingDQN}. Prioritized Experience Replay samples transitions according to learning signal rather than uniformly \parencite{Schaul2016PER}. Rainbow combined several improvements into a single DQN-family agent \parencite{Hessel2018Rainbow}.

Distributional reinforcement learning changes the object being estimated. Instead of only estimating expected return, it approximates the distribution of possible returns \parencite{Bellemare2017DistributionalRL}. C51 represents the distribution over a fixed categorical support, while QRDQN uses quantile regression to approximate return quantiles without predefining the same fixed support \parencite{Dabney2018QRDQN}. This makes QRDQN part of the value-based DQN family, but with a distributional target.

The security motivation is intuitive but must be stated carefully. False positives and false negatives have asymmetric consequences, and rare high-cost errors can matter more than a mean score suggests. A distributional value estimate may be useful for studying return structure under asymmetric rewards. It does not automatically provide calibrated uncertainty, operational safety, or risk-sensitive behavior. Those properties would require additional decision rules, calibration, or evaluation criteria beyond simply training a distributional agent \parencite{Bellemare2017DistributionalRL,Dabney2018QRDQN}.

In this thesis, QRDQN is an exploratory algorithmic choice. It is suitable for a discrete binary action space and moderately high-dimensional flow observations, and it is available through the SB3-Contrib implementation used by the project \parencite{SB3ContribQRDQNDocs,GymnasiumDocs}. It should not be described as proven superior for NIDS unless same-protocol comparisons against DQN, Random Forest, and other baselines support that claim.

\section{Reinforcement Learning for Intrusion Detection and Cyber Defense}

RL and DRL for intrusion detection already form an active literature \parencite{Yang2024DRLNIDSSurvey,Kheddar2025RLIDSReview,Gueriani2024DRLIoTIDSSurvey,Adawadkar2022CyberSecurityRLSurvey}. Prior work includes DQN-style intrusion classifiers, actor-critic approaches, adversarial-environment formulations, feature-selection agents, SDN or IoT-focused systems, and autonomous cyber-defense simulations. The field is heterogeneous, so papers should not be treated as directly comparable unless their task formulation, datasets, actions, rewards, and validation protocols are aligned.

One branch is close to this thesis: dataset-as-environment intrusion detection. Lopez-Martin et al. reformulated supervised intrusion detection as a deep RL problem by treating labeled records as states, class predictions as actions, and reward as a function of correctness \parencite{LopezMartin2020DRLIDS}. Their later work extended a related offline RL formulation across several datasets, including CICIDS2017 and UNSW-NB15 \parencite{LopezMartin2021RBFOfflineRL}. Ren et al. used DRL for feature selection and classification on CSE-CIC-IDS2018 \parencite{Ren2022IDRDRL}. These works show that RL-style formulations for IDS are precedented.

Another branch studies more genuinely sequential cyber defense. POMDP and attack-graph approaches model defensive actions under uncertainty \parencite{Hu2020POMDPDefense}. CybORG provides a gym-style environment for autonomous cyber agents \parencite{Standen2021CybORG}. Broader autonomous cyber defense surveys discuss agents that select countermeasures in simulated or emulated networks \parencite{Palmer2023_weak,CSET2023AutonomousCyberDefense}. These settings are conceptually different from this thesis because defender actions can alter future state.

The thesis therefore belongs primarily to the first branch: offline flow-level detection formulated through an RL interface. It can reference autonomous cyber defense as broader context and future direction, but it should not imply that the current system performs multi-agent, adversarial, or live network-control learning.

\section{Classification-as-RL and Binary Security Decisions}

The \texttt{PERMIT}/\texttt{BLOCK} formulation belongs in the classification-as-RL discussion. In the current environment, each flow observation is presented to the agent, the action is binary, and reward is computed from the action and ground-truth label. \texttt{PERMIT} corresponds to accepting traffic as benign, while \texttt{BLOCK} corresponds to identifying it as attack. A false negative permits an attack flow; a false positive blocks benign traffic.

This framing has a real methodological benefit: it makes cost-sensitive security decisions explicit. IDS errors are not symmetric in practice. Missing an attack and blocking benign traffic have different operational meanings, and the preferred trade-off depends on the defended environment \parencite{Lee2002CostSensitiveIDS,Cardenas2006IDSFramework}. Reward design can encode this asymmetry directly, which is useful for a thesis studying defensive decision rules rather than only label prediction.

The limitation is equally important. If the dataset does not react to the action, then \texttt{PERMIT}/\texttt{BLOCK} is not an active control action. It is an offline decision label over a recorded or extracted flow row. The formulation is therefore closer to reward-engineered classification than to a firewall, IPS, SDN controller, or autonomous cyber-defense agent. This is not a flaw if stated clearly. It becomes a flaw only if the chapter overclaims live blocking, deployment readiness, or fully sequential control.

Because the formulation is close to cost-sensitive classification, supervised baselines are mandatory. A strong Random Forest, and ideally other tabular baselines, help determine whether QRDQN adds empirical value beyond the canonical features, preprocessing, and reward design. The thesis should interpret any RL advantage as benchmark- and protocol-specific unless broader evidence is available.

\section{Methodological Pitfalls in ML/DL/RL-Based NIDS}

The NIDS literature contains repeated warnings against overinterpreting benchmark accuracy. Sommer and Paxson argued that ML-based intrusion detection is difficult to validate in realistic environments because benign traffic is diverse, attacks are rare, labels are hard to obtain, and operational base rates differ sharply from benchmark conditions \parencite{SommerPaxson2010ClosedWorld}. Arp et al. identify similar risks across machine learning in computer security, including data leakage, sampling bias, inappropriate metrics, and weak experimental design \parencite{Arp2020DosDontsMLSecurity}.

Leakage is broader than accidentally placing the same row in train and test. It includes any target-related information that would not legitimately be available at prediction time \parencite{Kaufman2011LeakageDataMining}. In NIDS, leakage can enter through endpoint identifiers, absolute timestamps, Flow IDs, scenario-specific ports, global preprocessing fitted before splitting, duplicate or near-duplicate flows, and train/test partitions that share capture context. This is especially important for CICIDS2017 because several later studies report errors or artifacts in flow extraction, labeling, and summarized CSV use \parencite{Engelen2021CICIDSIssues,Lanvin2023Errors,Liu2022ErrorPrevalenceNIDS,Rosay2022CICIDS2017Analysis,Dube2024FaultyCICIDS2017}.

Random row-wise splits are a common source of inflated confidence. They preserve class proportions and are useful for early internal checks, but they can allow related flows, day-specific patterns, or attack-scenario artifacts to appear on both sides of the split. Harder splits by day, capture file, scenario, or external dataset are stronger tests of generalization \parencite{Apruzzese2022CrossEvaluationNIDS,Layeghy2023CrossDomainNIDS,Cantone2024}.

RL introduces additional pitfalls. Deep RL training can be stochastic and sensitive to random seeds, hyperparameters, reward shaping, and environment details. In a static dataset-as-environment setting, strong results may reflect the same feature and split advantages that benefit supervised models. Therefore, any claim that QRDQN is better than a supervised baseline must be supported by same-protocol comparison and should acknowledge run-to-run variance when repeated seeds are not available.

\section{Evaluation Protocols, Metrics, and Reproducibility}

Evaluation should match the strength of the claim. Same-dataset random splits are acceptable for internal learning checks, but they are weak evidence for deployment. Day-based, file-based, leave-one-scenario-out, leave-one-exact-CSV-out, cross-dataset, and independently captured lab-traffic evaluations provide progressively stronger tests of distribution shift. This thesis follows that logic by separating CICIDS2017 internal benchmarking from stricter validation and offline Phase 2 inference on lab-derived flows.

Metrics must also be interpreted operationally. Accuracy is useful but insufficient for imbalanced NIDS. Precision indicates how trustworthy positive alerts are, recall indicates how much attack traffic is detected, and false-positive and false-negative rates map more directly to analyst workload and missed-threat risk \parencite{Axelsson1999BaseRate,Fawcett2006ROC,Saito2015PRROC}. Confusion matrices, per-class precision, recall, F1, TP, FP, FN, TN, FPR, and FNR should be preferred over a single headline number \parencite{Milenkoski2015IDSEvaluationSurvey}.

Cost-sensitive evaluation is relevant because the correct operating point depends on context. NIST guidance and IDS evaluation literature both emphasize trade-offs between false positives and false negatives \parencite{ScarfoneMell2007,Cardenas2006IDSFramework}. In this thesis, the asymmetric reward design should be treated as a scenario assumption and evaluated empirically, not as a universal cost model for all networks.

Reproducibility is part of the evidence, not an administrative detail. Cybersecurity ML results should preserve dataset version, feature schema, excluded fields, preprocessing order, split logic, scaler state, random seeds, model configuration, run ID, metrics, and predictions where feasible. Olszewski et al. show that reproducibility remains a major problem in ML security research \parencite{Olszewski2023ReproSecurityML}. The repository's use of saved configs, metrics, model artifacts, and validation outputs is therefore a methodological strength, but missing artifacts should remain clearly labeled as pending.

\section{Data Efficiency and Training-Scale Evaluation}

Data efficiency asks how much labeled traffic is needed to reach useful performance, and how the learning curve changes as the training set grows. This question is important for NIDS because labeled attack traffic is expensive, traffic distributions change, and public benchmark abundance does not imply deployment data abundance \parencite{Ring2019DatasetSurvey,Apruzzese2022CrossEvaluationNIDS}. It is also important for RL because deep RL methods can be sample-hungry, especially when reward signals are sparse, noisy, or highly shaped \parencite{SuttonBarto2018RL,Hessel2018Rainbow}.

Existing RL-IDS literature often trains on full public datasets and reports final benchmark scores, while controlled training-scale ablations are less common \parencite{Yang2024DRLNIDSSurvey,Kheddar2025RLIDSReview}. Supervised few-shot and class-incremental NIDS work addresses related questions from a different angle \parencite{DiMonda2024FewShotNIDS}, but the thesis question is narrower: under one feature schema and evaluation protocol, how does QRDQN scale with available training rows, and how does that compare with supervised baselines?

The thesis positions training-scale experiments as methodological evidence. Budgets such as 100k, 250k, 500k, 1M, and 2M samples can test whether the RL formulation requires substantially more data than a supervised baseline to reach comparable performance. Such experiments should not be framed as proof of few-shot learning or deployment readiness. They are internal evidence about the data requirements of this particular formulation.

\section{Research Gap and Positioning of This Thesis}

The literature supports a narrow but meaningful research gap. NIDS and flow-based ML are mature areas; public datasets such as CICIDS2017 are useful but fragile; supervised tabular baselines are strong; RL and DRL for IDS already exist; and autonomous cyber defense is a broader, more sequential field than the current implementation \parencite{Ring2019DatasetSurvey,LiuLang2019IDSSurvey,Yang2024DRLNIDSSurvey,Kheddar2025RLIDSReview}. The gap is not the absence of RL for intrusion detection.

The contribution of this thesis is a reproducible experimental study of an offline, flow-level, binary security-decision formulation. The system maps CICIDS2017 and later lab-flow inputs into a fixed 76-feature canonical schema, augments observations with a missingness mask to produce 152-dimensional inputs, exposes rows through a Gymnasium-compatible environment, and trains a QRDQN agent to choose between \texttt{PERMIT} and \texttt{BLOCK}. This makes the cost-sensitive decision interface explicit while preserving a clear comparison obligation against supervised baselines.

The thesis is also positioned methodologically. It treats CICIDS2017 as the main internal benchmark, not as proof of production readiness. It separates random-split results from stricter validation. It uses anti-leakage preprocessing and validation checks to reduce obvious artifact learning. It treats external lab-captured traffic as the preferred next evidential step, while acknowledging that current Phase 2 inference is offline and does not perform active blocking.

The defensible final claim is therefore not that QRDQN is universally better than supervised NIDS models, nor that the project implements autonomous cyber defense. The defensible claim is that the thesis evaluates a QRDQN-based, cost-sensitive, flow-level defender under a reproducible benchmark and validation pipeline, and uses supervised baselines, leakage-aware checks, training-scale analysis, and lab-traffic inference to determine how far that formulation can be trusted.
